% ============================================================
%  v4.2.0-r02 development patch
%  HS Split/Refold XOR Register Memory
%  Scope: source-only DEV-PATCH; no release claim, no PDF freeze.
% ============================================================

\subsection{HS split/refold XOR register memory}
\label{ssec:hs-split-refold-xor-register-memory}

\begin{remark}[Scope of the r02 split/refold register layer]
\label{rem:r02-split-refold-scope}
This subsection records a structural strengthening of the Hilbert-shift
ladder introduced in Subsection~\ref{ssec:op1-op4-joint-hcdixon-reduction}.
It is a DEV-PATCH layer only.  It does not claim a new global OP closure,
and it does not replace the OP1 compression proof or the finite OP4
selector proofs.  Its role is to explain how higher Hilbert shifts may
alternate between visible split and refolded center while preserving
route history in a reversible register.
\end{remark}

\begin{definition}[Even shift state]
\label{def:even-hs-state}
An even Hilbert-shift state is a pair
\[
  S_m=(K_m,R_m),
\]
where \(K_m\) is the effective refolded center at the \(m\)-th
split/refold cycle, and \(R_m\) is a finite register carrying the
accumulated route, lens, selector, orientation, and residual history.
The visible geometry is therefore allowed to remain small even when the
registered history grows.
\end{definition}

\begin{definition}[Focal split signature]
\label{def:focal-split-signature}
The next odd Hilbert-shift layer unfolds the effective center \(K_m\)
into a focal pair
\[
  K_m\longmapsto (K_m^-,K_m^+).
\]
The split produces a finite signature
\[
  \sigma_m.
\]
This signature may contain, depending on the active layer, branch data,
lens observables, selector data, residual data, and orientation data:
\[
  \sigma_m
  =
  (b_m,\Delta\theta_m,\Delta L_m,\chi_m,\rho_m,\omega_m).
\]
The entries are not required to be real-valued; they may be symbolic,
Boolean, selector-valued, or mask-valued, provided they live in the
declared finite register grammar.
\end{definition}

\begin{definition}[XOR-like register update]
\label{def:xor-like-register-update}
Let \((\mathcal X_m,\oplus)\) be the finite register algebra used at the
\(m\)-th split/refold step.  In the binary case \(\oplus\) is ordinary
XOR.  More generally, it is an explicitly declared reversible update
operation satisfying the inversion rule
\[
  (R_m\oplus\sigma_m)\oplus\sigma_m=R_m
\]
whenever the same signature \(\sigma_m\) is supplied.  The register update
is
\[
  R_{m+1}=R_m\oplus\sigma_m.
\]
Thus the update writes the shift signature without making the visible
refold irreversible.
\end{definition}

\begin{definition}[No-erase refold guard]
\label{def:no-erase-refold-guard}
A refold map
\[
  \operatorname{Refold}_m
\]
is admissible only if it compresses the visible focal pair
\[
  (K_m^-,K_m^+)
\]
to an effective center \(K_{m+1}\) without erasing the updated register
\(R_{m+1}\).  We record this as the guard
\[
  \boxed{\operatorname{Refold}\neq \operatorname{Erase}.}
\]
In the minimal midpoint model one writes
\[
  K_{m+1}
  =
  \frac12(K_m^-+K_m^+),
\]
but in the full register model the refold may depend on the written
history:
\[
  K_{m+1}
  =
  \operatorname{Refold}_m(K_m^-,K_m^+;R_{m+1}).
\]
\end{definition}

\begin{lemma}[HS split/refold recursion]
\label{lem:hs-split-refold-recursion}
Let \(S_m=(K_m,R_m)\) be an even Hilbert-shift state.  Suppose the next
odd shift produces a focal pair and signature
\[
  K_m\mapsto(K_m^-,K_m^+),
  \qquad
  \sigma_m,
\]
and suppose the register is updated by
\[
  R_{m+1}=R_m\oplus\sigma_m.
\]
If the following refold is no-erase admissible, then the next even state is
\[
  S_{m+1}
  =
  (K_{m+1},R_{m+1}),
\]
where
\[
  K_{m+1}
  =
  \operatorname{Refold}_m(K_m^-,K_m^+;R_{m+1}).
\]
The visible geometry may therefore alternate
\[
  1\longrightarrow2\longrightarrow1
\]
while the route history remains stored in the register chain.
\end{lemma}

\begin{remark}[Active geometry versus addressable history]
\label{rem:active-geometry-vs-history}
The split/refold recursion separates two notions which should not be
identified.  The active visible geometry is locally \(1\)- or \(2\)-centered:
\[
  K_m
  \quad\text{or}\quad
  (K_m^-,K_m^+).
\]
The addressable history space, however, may grow like a binary register
space.  For \(m\) binary signatures there are \(2^m\) possible histories,
but a single realized route carries only its register path.  Thus higher
Hilbert shifts need not produce a visible exponential point cloud; they
produce reversible exponential addressability.
\end{remark}

\begin{definition}[Higher-HS alternation convention]
\label{def:higher-hs-alternation-convention}
We read higher Hilbert-shift layers by the convention
\[
  \mathrm{HS}(2m): (K_m,R_m),
\]
and
\[
  \mathrm{HS}(2m+1): (K_m^-,K_m^+,R_m,\sigma_m).
\]
The subsequent even layer writes and refolds:
\[
  (K_m^-,K_m^+,R_m,\sigma_m)
  \longmapsto
  (K_{m+1},R_m\oplus\sigma_m).
\]
This convention refines the HS(0)-anchored ladder by treating odd shifts
as focal/split layers and even shifts as effective refold/write layers.
\end{definition}

\begin{remark}[Reading the low HS layers]
\label{rem:low-hs-layers-split-refold-reading}
Under this convention one may read the low layers as follows:
\[
\begin{array}{c|l}
\mathrm{HS}(0) & \text{raw forbidden kernel }K_0,\\
\mathrm{HS}(1) & \text{first focal/ellipsoid HC split }(K_0^-,K_0^+),\\
\mathrm{HS}(2) & \text{effective stable collar/orbit center }K_1,\\
\mathrm{HS}(3) & \text{Fano/octonionic split over }K_1,\\
\mathrm{HS}(4) & \text{selector/leakage refold over }K_2,\\
\mathrm{HS}(5) & \text{front/back or readout/action split over }K_2,\\
\mathrm{HS}(6) & \text{Dixon carrier refold }K_3.
\end{array}
\]
This table is a structural reading, not an additional claim that all layers
are globally solved.
\end{remark}

\begin{definition}[Lens signature of a split]
\label{def:lens-signature-of-split}
For a split layer with focal pair \((K_m^-,K_m^+)\), and two admissible
test readouts \(A,B\), define the lens signature schematically by
\[
  \mathcal L_m(A,B)
  =
  (\Delta\theta_m,\Delta L_m,\mu_m,\rho_m),
\]
where \(\Delta\theta_m\) records angular displacement, \(\Delta L_m\)
records path-length or time-delay displacement, \(\mu_m\) records
magnification or readout amplification, and \(\rho_m\) records residual
debt.  The split signature \(\sigma_m\) may include \(\mathcal L_m(A,B)\)
when the layer is read as a lens/readout layer.
\end{definition}

\begin{lemma}[XOR reconstruction under known signatures]
\label{lem:xor-reconstruction-known-signatures}
Assume that the signatures \(\sigma_0,\ldots,\sigma_{m-1}\) are stored or
deterministically reconstructible from the declared shift grammar.  If
\[
  R_{j+1}=R_j\oplus\sigma_j
\]
for each \(j<m\), then the register chain can be inverted stepwise:
\[
  R_j=R_{j+1}\oplus\sigma_j.
\]
Consequently, a refolded even state can preserve reconstructible route
history even though its visible geometry has been compressed to one
effective center.
\end{lemma}

\begin{lemma}[No higher-HS residual explosion criterion]
\label{lem:no-higher-hs-residual-explosion}
Suppose every higher-HS split produces a unique finite signature
\(\sigma_m\), every register update is reversible, and every refold
satisfies the no-erase guard.  Then higher Hilbert shifts do not create
an uncontrolled visible exponential residual geometry.  Any additional
history is routed into the finite register chain
\[
  R_m
  =
  R_0\oplus\sigma_0\oplus\cdots\oplus\sigma_{m-1}.
\]
Thus possible exponential addressability is represented as register
history rather than as uncontrolled new geometric branches.
\end{lemma}

\begin{theorem}[OP1/OP4 register-closure strengthening candidate]
\label{thm:op1-op4-register-closure-strengthening}
Assume the OP1/OP4 joint HC--Dixon reduction of
Theorem~\ref{thm:conditional-op1-op4-local-closure}.  Add the following
split/refold register assumptions:
\begin{enumerate}
  \item each pre-Dixon Hilbert-shift split writes a unique finite
  signature \(\sigma_m\);
  \item the register update is reversible,
  \(R_{m+1}=R_m\oplus\sigma_m\);
  \item the refold map satisfies \(\operatorname{Refold}\neq
  \operatorname{Erase}\);
  \item the signatures associated with the OP4 selectors are finite
  register data:
  \[
    \sigma_{15/16},\qquad
    \sigma_{4/3},\qquad
    \sigma_{7/6},\qquad
    \sigma_{\delta^2}.
  \]
\end{enumerate}
Then the OP1 compression route and the leading OP4 selector readouts may
be read as finite reversible register data over the OP5-inherited
HC/Dixon route grammar.  In particular, higher HS layers do not by
themselves create a new uncontrolled OP1/OP4 residual channel merely by
splitting, provided their split signatures are written and refolded
without erasure.
\end{theorem}

\begin{remark}[Status of the strengthening theorem]
\label{rem:register-strengthening-status}
Theorem~\ref{thm:op1-op4-register-closure-strengthening} is a
strengthening candidate, not a proof of the OP1 load identity or of the
finite OP4 selector values.  Its role is to explain how the pre-Dixon
Hilbert-shift ladder can remain geometrically controlled while still
carrying enough reversible memory to reconstruct route history and
selector origin.
\end{remark}

\begin{definition}[r02 proof-obligation register]
\label{def:r02-proof-obligation-register}
The r02 split/refold register layer has the following proof obligations:
\[
  \rho_{\mathrm{r02}}
  =
  (
    \rho_{\mathrm{sig}},
    \rho_{\oplus},
    \rho_{\mathrm{noerase}},
    \rho_{\mathrm{lens}},
    \rho_{\mathrm{selector\;trace}}
  ).
\]
Here \(\rho_{\mathrm{sig}}\) records signature uniqueness,
\(\rho_{\oplus}\) records reversible update validity,
\(\rho_{\mathrm{noerase}}\) records the no-erase refold guard,
\(\rho_{\mathrm{lens}}\) records the adequacy of lens signatures, and
\(\rho_{\mathrm{selector\;trace}}\) records the identification of
\(15/16,4/3,7/6,\delta_{\mathrm{CL}}^2\) as finite shift signatures.
\end{definition}

\begin{remark}[Impact on OP1 and OP4]
\label{rem:r02-impact-op1-op4}
The r02 register layer can help OP1 by giving a structural reason why
\(\gamma_L^5\) may be read as five controlled pre-Dixon signature stages
rather than as an uncontrolled geometric expansion.  It can help OP4 by
reading the finite selectors \(15/16\), \(4/3\), \(7/6\), and the
double-frame channel \(\delta_{\mathrm{CL}}^2\) as registered split/refold
signatures.  These are compatibility mechanisms, not substitutes for the
remaining selector, compression, or frame-extraction proofs.
\end{remark}
